% GAL1897 W02 -- Émile Picard, Introduction
% Assigned source: physical PDF pages 024--029 / JP2 leaves L0023--L0028.
% Logical roman folios v--x. The opening folio v is suppressed in the 1897 print;
% visible running folios VI--X are recorded at the exact source-page boundaries.
% UTF-8 LaTeX fragment. Requires graphicx for source-derived apparatus images and the terminal ornament.
% The PAGE_DISPOSITIONS.csv and EVIDENCE_MANIFEST.csv files are authoritative for
% physical-page coordinates, copy-specific marks, scan damage, and platform overlays.

\providecommand{\GalSourcePageStart}[4]{}%
\providecommand{\GalSourcePageBoundary}[4]{\space}%
\providecommand{\GalSourcePageBoundaryInWord}[4]{}%
\providecommand{\GalPrintedError}[2]{#1}%
\providecommand{\GalCopySpecificMarks}[4]{}%
\providecommand{\GalCopySpecificMarkImage}[4]{}%
\providecommand{\GalSourceDerivedOrnament}[1]{\includegraphics{#1}}%

\GalSourcePageStart{024}{L0023}{v}{suppressed}

% ENSEG:W02:PDF024:0001
\begin{center}
\rule{0.72\linewidth}{0.4pt}\par
\vspace{2.8em}
{\Large\bfseries INTRODUCTION.\par}
\vspace{2.2em}
\end{center}

% ENSEG:W02:PDF024:0002
The works of Galois were, as is well known, published%
\GalCopySpecificMarkImage{024}{L0023}{right-margin X-like copy-specific mark}{figures/PDF024_L0023_MARGIN_X_ENHANCED.png}%
\space in 1846 by Liouville in the \emph{Journal de Mathématiques}. It was regrettable
that the works of the great geometer could not be obtained as a separate volume;
the Mathematical Society therefore decided to have Galois's memoirs reprinted.
This edition conforms to the preceding one; only the notice placed by Liouville
at the beginning of the publication has been omitted.

% ENSEG:W02:PDF024:0003
A study of Galois's life, apparently definitive, has just been published by
M. Paul Dupuy in the \emph{Annales de l'École Normale supérieure} (1896).
Among the earlier documents concerning Galois's life, mention must be made of
the obituary notice devoted to him by his friend Auguste Chevalier in the
\emph{Revue encyclopédique} (September 1832), and of an article published in
the \emph{Magasin pittoresque} in 1848. Évariste Galois was born at
Bourg-la-Reine, near Paris, on 25 October 1811; he left his father's home in
1823 to enter the fourth form at the Collège Louis-le-Grand. From the age of
fifteen his extraordinary aptitude for the mathematical sciences began to
show itself. Elementary algebra books did not satisfy him, and he acquired
his algebraic education from Lagrange's classic works. It appears that by the
age of seventeen Galois had already obtained results of the highest importance
concerning the theory of algebraic equations. We can do no more than%
\GalSourcePageBoundary{025}{L0024}{vi}{VI}%
conjecture about the course of his ideas, since the two memoirs he submitted
to the Academy of Sciences were lost. One thing, however, is certain: at the
beginning of 1830 he possessed his general principles, as is shown by the
summary of a memoir on the algebraic solution of equations in the
\emph{Bulletin de Férussac}. It states a series of results that are plainly
applications of a general theory. This short article is the most important
one published by Galois himself, for the fundamental memoir on algebra found
among his papers was not printed until 1846.

% ENSEG:W02:PDF025:0004
M. Dupuy's article contains information of great interest about Galois's life.
It is unlikely that new documents will now be added to those we possess. After
twice failing the entrance examination for the École Polytechnique, Galois
entered the École Normale in 1829 and was compelled to leave it the following
year. During the last year of his life he devoted himself entirely to politics,
spent several months behind the bars of Sainte-Pélagie, and, mortally wounded
in a duel, died on 31 May 1832. Confronted with a life so short and so troubled,
our admiration redoubles for the prodigious genius who left so deep a mark on
science. Precocious achievement is not rare among great geometers, but Galois's
case is remarkable above all others. Alas, the unhappy young man seems to have
paid a bitter price for his genius. As his brilliant mathematical faculties
\GalPrintedError{develop}{W02-PE001}, his formerly cheerful and open character
darkens, and the consciousness of his immense superiority gives rise to
excessive pride. This caused the disappointments that so strongly affected his
career, the first being his failure at the École Polytechnique. His examination
at that%
\GalSourcePageBoundary{026}{L0025}{vii}{VII}%
school left vivid memories. Without going as far as legend would have us go,
let us say only that Galois refused to answer a question, which he considered
ridiculous, about the arithmetic theory of logarithms. Nor can there be any
doubt that he was unwilling to provide the explanations of his work requested
by the mathematicians with whom he came into contact---explanations made
necessary by the hurried composition of his memoirs. It is therefore easy to
understand why his merit was not recognized by his contemporaries. Liouville
did not succeed without difficulty in grasping the sequence of Galois's ideas,
and numerous commentators were still needed to fill gaps remaining in more
than one proof and to bring the great geometer's theories to the degree of
simplicity they are capable of assuming today.

\GalCriticalNote{W02-PE001}{\textbf{Status: proved typographical error.}
The 1897 French prints \emph{déveveloppent}; the intended reading is
\emph{développent}. The English sense is “develop.” This spelling correction
does not alter Picard's claim. The bounded post-P13 search found no explicit
prior notice securely dated in the searched corpus. Source: PDF025/L0024. Downstream effect: the English
spelling is repaired, but Picard's claim is unchanged. The primary 1897
reading and French morphology prove the duplication; no earlier independent
notice is established. Evidence: \texttt{W02-PE001} in
\nolinkurl{source/critical_baseline/PRINTED_AND_SOURCE_ERRORS.md} and
\texttt{BIBLIOGRAPHIC\_PRIOR\_NOTICE\_MATRIX.csv}.}

% ENSEG:W02:PDF026:0005
The theory of equations owes considerable advances to Lagrange, Gauss, and
Abel, but none of them succeeded in bringing to light the fundamental element
on which all the properties of an equation depend. That glory was reserved for
Galois, who showed that every algebraic equation has a corresponding group of
substitutions in which the equation's essential characteristics are reflected.
In algebra, group theory had previously been the subject of many investigations,
mostly by Cauchy, who had already introduced certain elements of classification.
Galois's study of the theory of equations showed him the importance of the
notion of an invariant subgroup of a given group. He was thus led to divide
groups into simple and composite groups, a fundamental distinction which in
fact extends far beyond algebra to the concept of groups of operations in its
broadest sense.

% ENSEG:W02:PDF026:0006
For general theories to acquire in science a%
\GalSourcePageBoundary{027}{L0026}{viii}{VIII}%
permanent right of citizenship, they usually need to be illustrated by
particular applications. In several fields such applications are not always
easy to find, and one could cite more than one theory in modern mathematics
confined, if I may say so, within its excessive generality. From the artistic
point of view, which plays a vital role in pure mathematics, nothing is more
satisfying than the development of these great theories; nevertheless,
thoughtful minds regret this tendency, which may have its dangers. No such
regret can be expressed in Galois's case. From the beginning, the algebraic
solution of equations provided him with a particular field of application in
which many geometers later followed him, foremost among them M. Camille Jordan.

% ENSEG:W02:PDF027:0007
Galois's work on algebraic equations made his name famous, but he appears to
have made discoveries in analysis that were at least as important. In his
letter to Auguste Chevalier, written on the eve of his death and constituting a
kind of scientific testament, Galois speaks of a memoir that could be composed
from his researches on integrals. We know of those researches only what he
says in this letter. Several points in some of Galois's statements remain
obscure; nevertheless, one can form an exact idea of some of the results he
had reached in the theory of integrals of algebraic functions. One thereby
becomes convinced that he possessed the most essential results concerning
Abelian integrals that Riemann was to obtain twenty-five years later. It is
with astonishment that we see Galois speak of the periods of an Abelian
integral associated with an arbitrary algebraic function. For hyperelliptic
integrals we have no difficulty understanding what he means by \emph{period},
but the general case is different, and%
\GalSourcePageBoundary{028}{L0027}{ix}{IX}%
one is almost tempted to suppose that Galois had at least glimpsed certain
notions concerning functions of a complex variable that were not to be
developed until several years after his death. The statements are precise:
the illustrious author classifies Abelian integrals into three kinds and
asserts that, if $n$ denotes the number of linearly independent integrals of
the first kind, there will be $2n$ periods. The theorem concerning inversion
of the parameter and argument in integrals of the third kind is clearly
indicated, as are the relations among the periods of Abelian integrals. Galois
also speaks of a generalization of Legendre's classical equation involving
the periods of elliptic integrals, a generalization that had probably led him
to the important relations since discovered by Weierstrass and M. Fuchs. We
have said enough to show the scope of Galois's discoveries in analysis. Had a
few more years been granted him to develop his ideas in this direction, he
would have been Abel's glorious successor and would have constructed, in its
essential parts, the theory of algebraic functions of one variable as we know
it today. Galois's reflections extended still farther. He ends his letter by
speaking of applying the theory of ambiguity to transcendental analysis. One
can roughly guess what he means, and in this field, which, as he says, is
immense, many discoveries still remain to be made even today.

% ENSEG:W02:PDF028:0008
It is not without emotion that one finishes reading the scientific testament
of this twenty-year-old, written on the eve of the day when he was to disappear
in an obscure quarrel. His death was an immense loss to science. Had Galois
lived, his influence would have greatly modi%
\GalSourcePageBoundaryInWord{029}{L0028}{x}{X}%
fied the direction of mathematical research in our country. I shall not risk
hazardous comparisons: Galois undoubtedly has equals among the great
mathematicians of this century; none surpasses him in the originality and
depth of his conceptions.

\begin{flushright}
\textsc{Émile Picard,}\\
{\small\bfseries President of the Mathematical Society of France.}
\end{flushright}

\begin{center}
\vspace{5.5em}
\GalSourceDerivedOrnament{figures/PDF029_L0028_END_ORNAMENT_ENHANCED.png}
\end{center}

\GalCopySpecificMarks{024--029}{L0023--L0028}{multiple non-letterpress marginal strokes and underlining}{Recorded page-by-page in EVIDENCE_MANIFEST.csv; not interpreted as Picard's printed text.}
